% GENERATED FILE. Edit canonical lessons, then run npm run generate:books.
% canonical-source-hash: fnv1a64:86b5720e79825b63
% canonical-lessons: GU-C108-vasan, GU-C108-chari, GU-C108-katar, GU-C108-soy, GU-C108-doro

\chapter{Vessel, Knife, Scissors, Needle, Thread}
\label{ch:gu-vessel-knife-scissors-needle-thread}
\hlchaptermodality{108}

\begin{chapteropening}
\textbf{By the end of this chapter:} \emph{I can say a vessel, a knife, scissors, a needle and a thread.}

Complete the last lesson of chapter 108: I can say a vessel, a knife, scissors, a needle and a thread.
\end{chapteropening}

\section[\texorpdfstring{\gu{વાસણ} (vāsaṇ)}{વાસણ (vāsaṇ)}]{\gu{વાસણ} (vāsaṇ) — a vessel, a utensil}

\label{lesson:GU-C108-vasan}



\begin{quote}
Before the new one: say the Gujarati for a step, then the Gujarati for a stove.
\end{quote}

\subsection*{You'll want to know: \gu{વાસણ}}
\textbf{\gu{વાસણ}} (\emph{vāsaṇ}) — \textquotedblleft{}a vessel, a utensil\textquotedblright{}.

\begin{cousinweb}
\textbf{\gu{વાસણ}} (\emph{vāsaṇ}) is a tadbhava, a word inherited from Sanskrit and reshaped by sound change. It comes from \emph{vāsana}, \textquotedblleft{}box, casket\textquotedblright{} (CDIAL 11599).
\end{cousinweb}

\begin{grammarlens}[title={What you've built}]
One more word for this chapter.
\end{grammarlens}

\subsection*{Your turn}
\begin{itemize}
\raggedright
  \item \emph{Say it:} \emph{vāsaṇ}
  \item \emph{Say it:} \emph{vāsaṇ}, once more
  \item \emph{Say it:} say \emph{vāsaṇ}
  \item \emph{Recall:} say the Gujarati for a step, then the Gujarati for a stove, then say \emph{vāsaṇ} again
\end{itemize}

\begin{tcolorbox}[breakable,colback=teal!4,colframe=teal!35!black,title={Before you move on}]
What does \gu{વાસણ} mean? (\textquotedblleft{}A vessel, a utensil\textquotedblright{}.) Say it once more.
\end{tcolorbox}
\section[\texorpdfstring{\gu{છરી} (charī)}{છરી (charī)}]{\gu{છરી} (charī) — a knife}

\label{lesson:GU-C108-chari}



\begin{quote}
Before the new one: say the Gujarati for a stove, then the Gujarati for a vessel.
\end{quote}

\subsection*{You'll want to know: \gu{છરી}}
\textbf{\gu{છરી}} (\emph{charī}) — \textquotedblleft{}a knife\textquotedblright{}.

\begin{cousinweb}
\textbf{\gu{છરી}} (\emph{charī}) is a tadbhava of Sanskrit \emph{kṣurī}, \textquotedblleft{}knife, dagger\textquotedblright{}, the feminine of \emph{kṣura}, \textquotedblleft{}razor\textquotedblright{} (CDIAL 3727). CDIAL rejects J. Bloch's derivation from Sanskrit \emph{tsaru}.
\end{cousinweb}

\begin{grammarlens}[title={What you've built}]
One more word for this chapter.
\end{grammarlens}

\subsection*{Your turn}
\begin{itemize}
\raggedright
  \item \emph{Say it:} \emph{charī}
  \item \emph{Say it:} \emph{charī}, once more
  \item \emph{Say it:} say \emph{charī}
  \item \emph{Recall:} say the Gujarati for a stove, then the Gujarati for a vessel, then say \emph{charī} again
\end{itemize}

\begin{tcolorbox}[breakable,colback=teal!4,colframe=teal!35!black,title={Before you move on}]
What does \gu{છરી} mean? (\textquotedblleft{}A knife\textquotedblright{}.) Say it once more.
\end{tcolorbox}
\section[\texorpdfstring{\gu{કાતર} (kātar)}{કાતર (kātar)}]{\gu{કાતર} (kātar) — scissors}

\label{lesson:GU-C108-katar}



\begin{quote}
Before the new one: say the Gujarati for a vessel, then the Gujarati for a knife.
\end{quote}

\subsection*{You'll want to know: \gu{કાતર}}
\textbf{\gu{કાતર}} (\emph{kātar}) — \textquotedblleft{}scissors\textquotedblright{}.

\begin{cousinweb}
\textbf{\gu{કાતર}} (\emph{kātar}) is a tadbhava, a word inherited from Sanskrit and reshaped by sound change. It comes from \emph{kartari}, \textquotedblleft{}scissors, knife\textquotedblright{} (CDIAL 2858).
\end{cousinweb}

\begin{grammarlens}[title={What you've built}]
One more word for this chapter.
\end{grammarlens}

\subsection*{Your turn}
\begin{itemize}
\raggedright
  \item \emph{Say it:} \emph{kātar}
  \item \emph{Say it:} \emph{kātar}, once more
  \item \emph{Say it:} say \emph{kātar}
  \item \emph{Recall:} say the Gujarati for a vessel, then the Gujarati for a knife, then say \emph{kātar} again
\end{itemize}

\begin{tcolorbox}[breakable,colback=teal!4,colframe=teal!35!black,title={Before you move on}]
What does \gu{કાતર} mean? (\textquotedblleft{}Scissors\textquotedblright{}.) Say it once more.
\end{tcolorbox}
\section[\texorpdfstring{\gu{સોય} (soy)}{સોય (soy)}]{\gu{સોય} (soy) — a needle}

\label{lesson:GU-C108-soy}



\begin{quote}
Before the new one: say the Gujarati for a knife, then the Gujarati for scissors.
\end{quote}

\subsection*{You'll want to know: \gu{સોય}}
\textbf{\gu{સોય}} (\emph{soy}) — \textquotedblleft{}a needle\textquotedblright{}.

\begin{cousinweb}
\textbf{\gu{સોય}} (\emph{soy}) is a tadbhava, a word inherited from Sanskrit and reshaped by sound change. It comes from \emph{sūcī}, \textquotedblleft{}needle\textquotedblright{} (CDIAL 13551).
\end{cousinweb}

\begin{grammarlens}[title={What you've built}]
One more word for this chapter.
\end{grammarlens}

\subsection*{Your turn}
\begin{itemize}
\raggedright
  \item \emph{Say it:} \emph{soy}
  \item \emph{Say it:} \emph{soy}, once more
  \item \emph{Say it:} say \emph{soy}
  \item \emph{Recall:} say the Gujarati for a knife, then the Gujarati for scissors, then say \emph{soy} again
\end{itemize}

\begin{tcolorbox}[breakable,colback=teal!4,colframe=teal!35!black,title={Before you move on}]
What does \gu{સોય} mean? (\textquotedblleft{}A needle\textquotedblright{}.) Say it once more.
\end{tcolorbox}
\section[\texorpdfstring{\gu{દોરો} (doro)}{દોરો (doro)}]{\gu{દોરો} (doro) — a thread}

\label{lesson:GU-C108-doro}



\begin{quote}
Before the new one: say the Gujarati for scissors, then the Gujarati for a needle.
\end{quote}

\subsection*{You'll want to know: \gu{દોરો}}
\textbf{\gu{દોરો}} (\emph{doro}) — \textquotedblleft{}a thread\textquotedblright{}.

\begin{cousinweb}
\textbf{\gu{દોરો}} (\emph{doro}) is a tadbhava, a word inherited from Sanskrit and reshaped by sound change. It comes from \emph{davara}, \textquotedblleft{}string\textquotedblright{} (CDIAL 6225).
\end{cousinweb}

\begin{grammarlens}[title={What you've built}]
One more word for this chapter.
\end{grammarlens}

\subsection*{Your turn}
\begin{itemize}
\raggedright
  \item \emph{Say it:} \emph{doro}
  \item \emph{Say it:} \emph{doro}, once more
  \item \emph{Say it:} say \emph{doro}
  \item \emph{Recall:} say the Gujarati for scissors, then the Gujarati for a needle, then say \emph{doro} again
\end{itemize}

\begin{tcolorbox}[breakable,colback=teal!4,colframe=teal!35!black,title={Before you move on}]
What does \gu{દોરો} mean? (\textquotedblleft{}A thread\textquotedblright{}.) Say it once more.
\end{tcolorbox}
